\section{Limitations}
\label{sec:limitations}

The present evidence comes from three development datasets, one corrected P1
run per dataset, and one real TRAIN/TEST split per dataset. These datasets do
not span the intended final benchmark, and 200 requested records per dataset
provide limited evidence about scaling. Repeated generation seeds, repeated
real-data splits, uncertainty intervals, and predeclared statistical
comparisons are pending.

Human evaluation is incomplete and modest in size: Reviewer A receives 120
records and Reviewer B overlaps on 40. Even after adjudication, specialist
domain validity may be difficult for general annotators to judge, particularly
when feature names encode scientific measurements without sufficient domain
context. The final paper must report reviewer qualifications and any domain
guidance without exposing personal information.

Generation and B3 judging use the same local model. Their errors may therefore
be correlated, and findings may not generalize to other local models, hosted
models, prompts, decoding configurations, or model scales. B3's score is
uncalibrated. The proposed calibration, evidence fusion, and risk controller
have not been implemented or tested, so the current system provides no
calibrated risk control. Repair and richer dataset optimization are
unimplemented extensions.

Extreme imbalance makes single-split utility unstable. Ecoli4 and Abalone19
have only four and six positive records, respectively, in the committed TEST
splits. Several logistic-regression TSTR outputs are undefined, and the Ecoli4
generation pool contains 201 records despite a request for 200. Final utility
claims require correction or explicit accommodation of these issues.

Copying and nearest-neighbour support are diagnostics, not privacy measures.
Support distance depends on standardization, categorical encoding, feature
weighting, and dimensionality; it has no calibrated threshold in the current
study. Privacy and fairness validators, near-copy analysis, and subgroup-risk
evaluation are not implemented. Exact absence of a match to an unshown TRAIN
record must not be read as a privacy guarantee.

The historical P1 manifests record a dirty working tree, although they retain
resolved configuration, commit identity, prompts, traces, and outputs. The
implementation workbook could not be inspected with the required read-only
spreadsheet runtime during manuscript initialization and is not used as
evidence. Finally, the literature review is substantial but incomplete;
several repository notes remain unverified placeholders and are deliberately
excluded from the bibliography.
